\chapter{That nourishes human hope (My Very Best Friend)}
In the final episode of Madoka Magica, failure is victory is loss is triumph.

We open with the same tableau that ended the previous episode. Four figures remain, the key players in this apocalypse, for apocalypse it is: Every timeline we have seen has ended with a fight against Walpurgisnacht. There is nothing beyond the fight with her, because Homura keeps resetting the universe before the future can occur. Homura is the first figure, broken and bleeding, the sad clown who is endlessly victimized by her desperate attempts to find meaning in an absurd and uncaring universe. Laughing at her mockingly is the instrument of her defeat, Walpurgisnacht, the Harlequin who signifies that absurdity. Between them is Kyubey, the director, author, orchestrator, the master manipulator who choreographs their dance to please his unseen audience and thus derive power and sustenance from their emotional arcs.

But then there is Madoka. She has been inert, the prize the others fight over, but now at last she makes her choice. And what a choice it is: death. She will become death, the destroyer of worlds, slaying all witches at the moment of their birth, until ultimately there are none left but she herself, and then she will kill herself. 

But Madoka is not Sayaka. This is not suicide; this is the transcendent death, the death of the ego that gives access to eternity and unity \cite{ref87}. There is no more Madoka; she is an existence without beginning or end, and within her all things are one. 

First, however, a long-promised cake. That was the agreement between her and Mami, after all: that if Madoka could find nothing to wish for, they would share a cake. But Madoka just made her wish; why, therefore, does she have cake with Mami?

Because Madoka has solved the paradox. To become enlightened, to escape the karmic cycle of hope and despair in which the magical girls are trapped, one must shed all desire. But if one sheds all desire, including the desire to transcend, why would anyone transcend? Madoka has found the answer: the death of ego, the erasure of the self-other distinction, which eliminates desire because the subject doing the desiring and the object of the desire are one and the same. "If you meet a Buddha, kill the Buddha." \cite{ref88} Though she made the wish only moments ago (insofar as that concept can mean anything now that she exists equally throughout all of time), she no longer has any wishes, so she receives enlightenment and cake. 

There is more to this scene than cake, however. Madoka is handed back her notebook of costume designs by Mami, the protector and signifier of the traditional magical girl show. When Mami died, Madoka stopped talking about becoming a magical girl as a way to find purpose; instead, it became a sacrifice she repeatedly considered making for the good of others. By returning the notebook, Mami is symbolically passing the role of guardian of the magical girl tradition to Madoka, while at the same time restoring the idea that being a magical girl can be a calling rather than a sacrifice. 

Because if there is one thing Madoka most definitely is not, it's a martyr. She is not a Christ-figure, suffering and dying as a way of absorbing the sins of others; she explicitly destroys the witch-aspect of herself which carries that suffering. She is egoless and transcendent, and thus cannot suffer. Her role in this Faustian tale is that of Gretchen \cite{ref89}, and as such she is more of a Marian figure, pure and unsullied, interceding to obtain a state of grace for others. Except even that is not quite accurate, because Madoka doesn't intercede or plead mercy for the magical girls. They still become witches and die; the only change is that their witch-forms do not exist in this world, because Madoka erases them at the moment of creation. They still suffer and still despair, still die-but such is the nature of living in this world. Madoka's role is as guide and teacher, a psychopomp who carries the magical girls out of the world before they can become a problem for it. With her in her pure land \cite{ref90}, they learn, and perhaps someday transcend as she has. Meanwhile, on Earth, things are imperfect, but better. 

"Daijobu." "It will be all right." This is what Madoka tells Homura just before she appears in the final form variously dubbed Madokami or Godoka by fans (though Madokanon \cite{ref91} would be more appropriate, as she is more Bodhisattva than divinity). It is a powerful phrase in the iconographic roots of the show; in Cardcaptor Sakura it was the ultimate spell the heroine created at the end of the series, an expression of hope of nigh-limitless power \cite{ref92}. Madoka is already carrying out her duties as the warden and guardian and magical girls past. At the same time, however, her transformation sequence is brief, nonsexualized, and strongly implies her costume to be made from an Anthony-the familiars that dominated the witch's labyrinth in the first episode, the first instance of the strange and wild new aesthetic the show introduced.

Madoka is becoming a bridge between the old genre and the new. She speaks the assurances of the old genre to the representative of the new one. She gives her ribbons-chosen for her by Junko, who has repeatedly been paralleled to Mami-to Homura as well. It is not a complete restoration of the magical girl tradition-the new world is still dark, and being a magical girl is fraught with dangers and likely to end with death-but a partial restoration, acknowledging that there were good stories, good characters, and true themes to be found among magical girl shows past.

Chief among those themes is hope. Naïve hope, the optimistic belief that things will get better, is a trap, yes. Anyone sitting around and waiting for a savior or a lucky break is doomed to disappointment. It is the nature of an entropic universe that if things can get worse, they will, and things can always get worse. But there is another form of hope, the hope embraced in the end by Homura. If things can and will get worse, then that necessarily means that at this moment, the universe is not at maximum awfulness; there must be something good in the world right now. That good can be sought out. It can be fought for, preserved for a little while. Entropy can be reversed locally. 

Madoka has attained enlightenment and divorced herself from this decaying world. But she has not abandoned it; the world she creates is better. Not perfect, because a perfect world is a world devoid of story, but better. The magical girls still inevitably die, but so does everyone else; what's important is that they now have a far better idea of how the system works and a much better relationship with each other and the Incubators-notably, the fact that wraiths drop a number of little magic-restoratives rather than one big one encourages the magical girls to work together. Teams are likely the norm in this new world, rather than solitary girls as in the old world, and since the Incubators can no longer derive energy from the despair of the witches, they have no incentive to make the girls suffer or hide from them how the system works.

Even Junko is shown in a new environment. We have seen her driven and determined before, concerned, caring, but this final sequence is the first time that we see her being happy. It's conceivable that this scene implies that Junko is a very different person in the new timeline, less driven and more nostalgic, but there's little reason to believe this is the case. It seems highly unlikely Madoka would replace her mother with a different woman, and far more likely that this is what Junko is like when she's relaxing and having a day out with her family. Her dynamic with Madoka's father is unchanged-he cares for Tetsuya while Junko deals with the outside world, in this case talking to Homura-and so it is likely that she is still the primary earner, the driven executive. It is simply that we can now see that she also contains within herself nostalgia and serenity and wistfulness; she contains contradictions, just as the magical girls/witches contain both curses and blessings, as this ending is both happy and sad, a win and a loss.

Seeing Junko and Tetsuya helps Homura to understand that there can be good things in what for her is a dark, Madoka-less world. She continues on, affirmed in her knowledge that Madoka is all around her, even if she cannot see her. She does not fight for hope in the normal sense, but out of love, and duty, and hope in the Havelian sense that whether or not she succeeds, her life makes sense as long as she fights. And so she fails to save Madoka, and in her failure succeeds in empowering Madoka to save herself. Madoka saves herself by sacrificing herself, and Homura loses her-but someday, when Homura expends the last of her energy and loses her last battle as a magical girl, she will be together with Madoka again.

But this is not for Homura alone. Someone else has been working, trying to stave off decay, but increasingly concerned that their efforts are doomed. "In order to write a perfect ending for a story you have to twist the laws of cause and effect," Urobuchi wrote in the afterword we have quoted twice before, "reverse black and white, and even possess a power to move in the opposite direction from the rule of the universe." \cite{ref93} The implied author of that note and this series is a deeply depressed individual, spiraling into a creative abyss brought on by despair.

"Only a heavenly and chaste soul that can sing carols of praise towards humanity can save the story." And now, in Madoka, all things are one. This is fiction, a creation within the mind of an author (even the gestalt implied author of a collaboration); the author is that one. Madoka loved something in the world enough to deem it worth saving, and she is part of that author. Homura will accept that love as reason enough to keep moving and working, and she is part of that author. Just as Homura is not suddenly all smiles and laughs in the new world, this is not a panacea-but it is enough to keep going for a while longer.

-Don't forget. Always, somewhere, someone is fighting for you.

-As long as you remember her, you are not alone.

The projector winds to a stop.
